\begin{document}

\begin{center}
{\Large\textbf{Основание, высота и плоское сечение}}\\[2mm]
{\small ЕГЭ, стереометрия, задание №3 \quad | \quad Фамилия\ \answerline{6cm}}
\end{center}

\ifteacher
\begin{center}\fbox{\textbf{Учительская версия:} ответы и краткие опоры выделены курсивом.}\end{center}
\fi

\textbf{Цель.} Перед вычислением назови тело, его основание и высоту либо
выбери плоское сечение, в котором находятся нужные отрезки. Рисунок является
схемой: измерять на нём ничего не нужно.

\section*{1. Осевое сечение конуса}

\stereotask{1}{Высота конуса равна \(18\), а длина образующей равна \(30\).
Найдите площадь осевого сечения этого конуса.}{\csname figStereoV05\endcsname}
{В осевом сечении \(r=\sqrt{30^2-18^2}=24\). Площадь равна \(rh=432\).}

\stereotask{2}{Диаметр основания конуса равен \(32\), а длина образующей равна
\(20\). Найдите площадь осевого сечения этого конуса.}{\csname figStereoV06\endcsname}
{Радиус \(16\), высота \(\sqrt{20^2-16^2}=12\); площадь сечения \(rh=192\).}

\section*{2. Призма и цилиндр}

\stereotask{3}{Цилиндр вписан в правильную четырёхугольную призму. Радиус
основания и высота цилиндра равны \(3\). Найдите площадь боковой поверхности
призмы.}{\csname figStereoV01\endcsname}
{Сторона квадрата \(6\), периметр \(24\), высота \(3\). Ответ: \(72\).}

\stereotask{4}{Цилиндр вписан в правильную шестиугольную призму. Радиус
основания цилиндра равен \(\sqrt3\), а высота призмы равна \(2\). Найдите
площадь боковой поверхности призмы.}{\csname figStereoV02\endcsname}
{Сторона правильного шестиугольника равна \(2\), периметр \(12\). Ответ:
\(12\cdot2=24\).}

\stereotask{5}{В основании прямой призмы лежит прямоугольный треугольник с
катетами \(5\) и \(6\). Боковые рёбра призмы равны \(\dfrac4\pi\). Найдите
объём цилиндра, описанного около этой призмы.}{\csname figStereoV36\endcsname}
{Гипотенуза \(\sqrt{61}\), радиус описанной окружности
\(\sqrt{61}/2\). \(V=\pi\cdot61/4\cdot4/\pi=61\).}

\section*{3. Диагональное сечение и угол между прямыми}

\stereotask{6}{В правильной четырёхугольной пирамиде \(SABCD\) точка \(O\) —
центр основания, \(SO=9\), \(SC=15\). Найдите длину отрезка \(BD\).}
{\csname figStereoV31\endcsname}
{\(OC=\sqrt{15^2-9^2}=12\), поэтому \(BD=2OC=24\).}

\stereotask{7}{В кубе \(ABCDA_1B_1C_1D_1\) найдите угол между прямыми
\(DC_1\) и \(BD\). Ответ дайте в градусах.}{\csname figStereoV11\endcsname}
{Для направляющих векторов скалярное произведение по модулю равно половине
произведения длин, поэтому угол между прямыми \(60^\circ\).}

\stereotask{8}{В правильной треугольной призме \(ABCA_1B_1C_1\), все рёбра
которой равны \(2\), найдите угол между прямыми \(BB_1\) и \(AC_1\). Ответ
дайте в градусах.}{\csname figStereoV12\endcsname}
{\(AA_1\parallel BB_1\). В прямоугольном треугольнике \(AA_1C_1\)
\(AA_1=AC=2\), поэтому угол равен \(45^\circ\).}

\ifteacher
\section*{Короткая проверка урока}
Ученик готов перейти к объёмам частей многогранника, если перед формулой
указывает основание и высоту, а для Пифагора называет конкретное сечение.
\fi

\end{document}
